\chapter{Cronograma del proyecto}
\label{anexo:cronograma}

Este anexo recoge la planificación temporal seguida durante el desarrollo de TierraViva, desde el inicio del Trabajo de Fin de Grado a finales de agosto de 2025 hasta la entrega y defensa previstas en julio de 2026. El cronograma no se plantea como una secuencia rígida y lineal, sino como la reconstrucción ordenada del trabajo real realizado: definición inicial, revisión tecnológica, primeros prototipos, reajuste de arquitectura, implementación \textit{hardware} y \textit{software}, validación funcional, prueba en Agrolab Móstoles y redacción progresiva de la memoria.

TierraViva ha requerido integrar electrónica, \textit{firmware}, comunicaciones móviles, plataforma \textit{cloud}, base de datos, API, \textit{dashboard}, actuación física sobre bomba de agua y documentación técnica. Por ello, varias tareas se han desarrollado en paralelo y otras han requerido iteraciones, correcciones y cambios de enfoque.

\section{Periodo de desarrollo}

El desarrollo efectivo del proyecto se extiende desde finales de agosto de 2025 hasta julio de 2026. La entrega de la memoria se sitúa entre los días 6 y 11 de julio de 2026, y la defensa está prevista entre los días 21 y 24 de julio de 2026.

\small
\renewcommand{\arraystretch}{1.15}
\rowcolors{2}{TableRowSoft}{white}

\begin{longtable}{|p{0.22\textwidth}|p{0.64\textwidth}|}
\caption{Resumen temporal del desarrollo del proyecto TierraViva.}
\label{cuadro:resumen_temporal_cronograma}\\
\hline
\rowcolor{URJCRed}
\TVHead{Periodo} & \TVHead{Caracterización del trabajo realizado} \\
\hline
\endfirsthead

\hline
\rowcolor{URJCRed}
\TVHead{Periodo} & \TVHead{Caracterización del trabajo realizado} \\
\hline
\endhead

Finales de agosto--octubre de 2025 & Inicio del TFG, definición de la idea, planificación inicial, primera propuesta al tutor y revisión bibliográfica. Septiembre y octubre tuvieron una dedicación más orientada a análisis, búsqueda de información y organización del alcance. \\
\hline
Noviembre--diciembre de 2025 & Consolidación del estado del arte, definición inicial de requisitos y arquitectura, primeros enfoques de prototipado local y selección inicial de tecnologías. \\
\hline
Enero--febrero de 2026 & Trabajo con Arduino, sensores y primeros montajes. Reajuste progresivo desde un enfoque más local hacia una arquitectura con estaciones remotas, conectividad LTE y ThingSpeak. Enero tuvo menor disponibilidad por exámenes. \\
\hline
Marzo--abril de 2026 & Desarrollo intensivo de comunicaciones, backend, base de datos, API y soporte multiestación. En abril se realizó una visita previa a Agrolab Móstoles junto con Ana Zazo, co-tutora del proyecto. \\
\hline
Mayo de 2026 & Consolidación del \textit{firmware} de estación, decisión local de riego, \textit{cooldown}, publicación HTTP, redacción seria de la memoria y cierre de los primeros capítulos técnicos. Mayo tuvo menor disponibilidad por exámenes, compensada con jornadas intensivas fuera del periodo de evaluación. \\
\hline
Junio de 2026 & Montaje de la estación B, versión final del \textit{dashboard}, ajustes eléctricos, validación funcional, prueba en Agrolab Móstoles el 22 de junio, fotografías, vídeo, evidencias y cierre de capítulos finales y anexos. \\
\hline
Julio de 2026 & Revisión final de memoria, formato, trazabilidad, entrega y preparación de la defensa. \\
\hline
\end{longtable}

\normalsize

\section{Fases principales del proyecto}

La Tabla~\ref{cuadro:fases_cronograma} resume las fases principales del desarrollo. Se ha separado la planificación inicial del reajuste posterior de arquitectura porque este cambio fue relevante para el resultado final del prototipo.

\small
\renewcommand{\arraystretch}{1.17}
\rowcolors{2}{TableRowSoft}{white}

\begin{longtable}{|p{0.12\textwidth}|p{0.20\textwidth}|p{0.56\textwidth}|}
\caption{Fases principales del desarrollo de TierraViva.}
\label{cuadro:fases_cronograma}\\
\hline
\rowcolor{URJCRed}
\TVHead{Fase} & \TVHead{Periodo} & \TVHead{Descripción} \\
\hline
\endfirsthead

\hline
\rowcolor{URJCRed}
\TVHead{Fase} & \TVHead{Periodo} & \TVHead{Descripción} \\
\hline
\endhead

F1 & Ago.--sep. 2025 & Inicio del TFG, análisis del problema, motivación del proyecto y preparación de la primera propuesta formal al tutor. La idea se orientó desde el inicio hacia un sistema IoT aplicado al riego inteligente y a la supervisión ambiental. \\
\hline
F2 & Sep.--nov. 2025 & Revisión bibliográfica y estado del arte: IoT, agricultura de precisión, riego inteligente, redes de sensores, plataformas embebidas, conectividad en campo, protocolos de aplicación y plataformas \textit{cloud}. \\
\hline
F3 & Sep.--dic. 2025 & Definición inicial de requisitos, arquitectura preliminar y alcance. Se estudiaron alternativas de sensorización, actuación, supervisión y comunicación, incluyendo enfoques locales basados en Raspberry Pi, puerto serie y MQTT. \\
\hline
F4 & Dic. 2025--feb. 2026 & Primeras pruebas de prototipado: Arduino UNO compatible, sensores BME280, BH1750 y sensor capacitivo de humedad de suelo. Resolución de problemas de programación asociados al conversor USB-serie CH340G. \\
\hline
F5 & Ene.--feb. 2026 & Reajuste progresivo de arquitectura hacia estaciones remotas con decisión local, conectividad LTE mediante A7670E-LASA y ThingSpeak como plataforma intermedia de telemetría. \\
\hline
F6 & Mar.--abr. 2026 & Trabajo con A7670E-LASA, comandos AT, conectividad LTE, publicación HTTP, ThingSpeak, base de datos SQLite, API FastAPI, soporte multiestación y backend en Raspberry Pi. \\
\hline
F7 & Mayo 2026 & Consolidación del \textit{firmware} de estación: lectura de sensores, modo seguro, confirmación de sequedad, \textit{cooldown}, decisión local de riego, actuación temporizada y publicación de estado. \\
\hline
F8 & Mayo--jun. 2026 & Redacción intensiva de memoria en paralelo al cierre técnico: capítulos de diseño, implementación, software, anexos técnicos, plan de pruebas, riesgos, trazabilidad y manuales. \\
\hline
F9 & Jun. 2026 & Montaje y prueba de la estación B, versión final del \textit{dashboard}, ajustes del bloque de bomba y alimentación, incorporación del condensador de \(2200~\mu\mathrm{F}\), capturas, logs y evidencias. \\
\hline
F10 & Jun. 2026 & Prueba funcional en Agrolab Móstoles. Hubo una visita previa en abril y una prueba final el 22 de junio, en la que se validó la estación B en entorno real de cultivo, se grabó vídeo y se recogió retroalimentación cualitativa. \\
\hline
F11 & Jun.--jul. 2026 & Cierre de resultados, discusión, conclusiones, revisión de coherencia, formato, anexos, entrega de la memoria y preparación de la defensa. \\
\hline
\end{longtable}

\normalsize

\section{Cronograma general}

La Tabla~\ref{cuadro:cronograma_general} recoge la distribución temporal de las actividades principales. Se utiliza una marca \textbf{X} para indicar los meses en los que cada actividad tuvo una dedicación significativa. Algunas tareas aparecen en varios meses porque se desarrollaron de forma iterativa.

\begin{landscape}
\scriptsize
\renewcommand{\arraystretch}{1.18}
\setlength{\tabcolsep}{3pt}
\rowcolors{2}{TableRowSoft}{white}

\begin{longtable}{|p{0.31\linewidth}|c|c|c|c|c|c|c|c|c|c|c|c|}
\caption{Cronograma general del proyecto TierraViva entre agosto de 2025 y julio de 2026.}
\label{cuadro:cronograma_general}\\
\hline
\rowcolor{URJCRed}
\TVHead{Actividad} &
\TVHead{Ago. 25} & \TVHead{Sep. 25} & \TVHead{Oct. 25} & \TVHead{Nov. 25} &
\TVHead{Dic. 25} & \TVHead{Ene. 26} & \TVHead{Feb. 26} & \TVHead{Mar. 26} &
\TVHead{Abr. 26} & \TVHead{May. 26} & \TVHead{Jun. 26} & \TVHead{Jul. 26} \\
\hline
\endfirsthead

\hline
\rowcolor{URJCRed}
\TVHead{Actividad} &
\TVHead{Ago. 25} & \TVHead{Sep. 25} & \TVHead{Oct. 25} & \TVHead{Nov. 25} &
\TVHead{Dic. 25} & \TVHead{Ene. 26} & \TVHead{Feb. 26} & \TVHead{Mar. 26} &
\TVHead{Abr. 26} & \TVHead{May. 26} & \TVHead{Jun. 26} & \TVHead{Jul. 26} \\
\hline
\endhead

Inicio del TFG, motivación y definición general de la idea & X & X &  &  &  &  &  &  &  &  &  &  \\
\hline
Primera propuesta formal al tutor &  & X &  &  &  &  &  &  &  &  &  &  \\
\hline
Revisión bibliográfica y estado del arte &  & X & X & X &  &  &  &  &  &  &  &  \\
\hline
Definición inicial de requisitos y arquitectura &  & X & X & X & X &  &  &  &  &  &  &  \\
\hline
Análisis de alternativas de arquitectura IoT &  & X & X & X & X & X & X &  &  &  &  &  \\
\hline
Primer enfoque local con Raspberry, puerto serie y MQTT &  &  &  & X & X & X &  &  &  &  &  &  \\
\hline
Reajuste hacia estaciones LTE y ThingSpeak &  &  &  &  &  & X & X & X & X &  &  &  \\
\hline
Selección y adquisición inicial de componentes &  &  &  & X & X & X & X &  &  &  &  &  \\
\hline
Compra de A7670E, relé, regulador y baterías &  &  &  &  &  &  & X &  &  &  &  &  \\
\hline
Compra de bomba y depósito de agua &  &  &  &  &  &  &  &  &  &  & X &  \\
\hline
Preparación de Arduino UNO y resolución CH340G &  &  &  &  &  & X &  &  &  &  &  &  \\
\hline
Validación de sensores BME280, BH1750 y humedad de suelo &  &  &  &  & X & X &  &  &  &  &  &  \\
\hline
Montaje físico inicial de la estación A &  &  &  &  &  & X & X &  &  &  &  &  \\
\hline
Trabajo con A7670E-LASA y comandos AT &  &  &  &  &  &  &  & X & X &  &  &  \\
\hline
Primera publicación correcta en ThingSpeak &  &  &  &  &  &  &  &  & X &  &  &  \\
\hline
Backend en Raspberry Pi &  &  &  &  &  &  &  & X & X & X &  &  \\
\hline
Diseño de base de datos SQLite &  &  &  &  &  &  &  & X & X &  &  &  \\
\hline
Implementación de API FastAPI &  &  &  &  &  &  &  & X & X &  &  &  \\
\hline
Soporte multiestación en backend y \textit{dashboard} &  &  &  &  &  &  &  & X & X &  & X &  \\
\hline
Primera demo básica de \textit{dashboard} &  &  &  &  &  & X &  &  &  &  &  &  \\
\hline
Segunda versión del \textit{dashboard} &  &  &  &  &  &  &  &  &  & X &  &  \\
\hline
Versión final del \textit{dashboard} &  &  &  &  &  &  &  &  &  &  & X &  \\
\hline
Desarrollo del \textit{firmware} estable de estación &  &  &  &  &  &  &  & X & X & X &  &  \\
\hline
Decisión local de riego en \textit{firmware} &  &  &  &  &  &  &  &  &  & X & X &  \\
\hline
Actuación con relé, bomba, temporización y \textit{cooldown} &  &  &  &  &  &  &  &  &  & X & X &  \\
\hline
Logs de monitor serie de estación &  &  &  &  &  &  &  &  &  & X &  &  \\
\hline
Montaje de estación B &  &  &  &  &  &  &  &  &  &  & X &  \\
\hline
Pruebas domésticas/laboratorio de estación B &  &  &  &  &  &  &  &  &  &  & X &  \\
\hline
Ajustes eléctricos y condensador de \(2200~\mu\mathrm{F}\) &  &  &  &  &  &  &  &  &  &  & X &  \\
\hline
Visita previa a Agrolab Móstoles &  &  &  &  &  &  &  &  & X &  &  &  \\
\hline
Prueba funcional en Agrolab Móstoles &  &  &  &  &  &  &  &  &  &  & X &  \\
\hline
Grabación de vídeo, fotografías y capturas finales &  &  &  &  &  &  &  &  &  &  & X &  \\
\hline
Redacción seria de memoria &  &  &  &  &  &  &  &  &  & X & X & X \\
\hline
Redacción de capítulos 1--6 &  &  &  &  &  &  &  &  &  & X & X &  \\
\hline
Elaboración de anexos técnicos &  &  &  &  &  &  &  &  &  & X & X & X \\
\hline
Redacción de capítulos 7 y 8 &  &  &  &  &  &  &  &  &  &  & X &  \\
\hline
Revisión final, formato, entrega y defensa &  &  &  &  &  &  &  &  &  &  & X & X \\
\hline

\end{longtable}
\end{landscape}

\section{Hitos principales}

La Tabla~\ref{cuadro:hitos_cronograma} identifica los hitos más relevantes del proyecto. Estos hitos permiten entender la evolución técnica de TierraViva y muestran que el resultado final no fue una implementación directa desde una solución cerrada, sino una construcción progresiva con reajustes.

\small
\renewcommand{\arraystretch}{1.16}
\rowcolors{2}{TableRowSoft}{white}

\begin{longtable}{|p{0.12\textwidth}|p{0.22\textwidth}|p{0.54\textwidth}|}
\caption{Hitos principales alcanzados durante el desarrollo de TierraViva.}
\label{cuadro:hitos_cronograma}\\
\hline
\rowcolor{URJCRed}
\TVHead{Hito} & \TVHead{Periodo} & \TVHead{Resultado alcanzado} \\
\hline
\endfirsthead

\hline
\rowcolor{URJCRed}
\TVHead{Hito} & \TVHead{Periodo} & \TVHead{Resultado alcanzado} \\
\hline
\endhead

H1 & Ago.--sep. 2025 & Inicio del TFG y definición de TierraViva como sistema IoT aplicado a riego inteligente. \\
\hline
H2 & Sep. 2025 & Entrega de la primera propuesta formal al tutor. \\
\hline
H3 & Sep.--nov. 2025 & Revisión bibliográfica y construcción del estado del arte sobre IoT, agricultura de precisión, riego inteligente y conectividad en campo. \\
\hline
H4 & Nov. 2025--ene. 2026 & Primer enfoque local con Raspberry Pi, puerto serie y MQTT, útil como aprendizaje técnico pero posteriormente reajustado. \\
\hline
H5 & Ene.--feb. 2026 & Decisión de avanzar hacia estaciones remotas con conectividad LTE, ThingSpeak y mayor autonomía del nodo de campo. \\
\hline
H6 & Ene. 2026 & Preparación de Arduino UNO compatible, instalación del controlador CH340G y validación de programación. \\
\hline
H7 & Dic. 2025--ene. 2026 & Validación inicial de sensores BME280, BH1750/GY-302 y sensor capacitivo de humedad de suelo. \\
\hline
H8 & Abr. 2026 & Primera publicación correcta en ThingSpeak mediante A7670E-LASA y HTTP. \\
\hline
H9 & Mar.--abr. 2026 & Desarrollo de backend en Raspberry Pi, base de datos SQLite, API FastAPI y soporte multiestación. \\
\hline
H10 & Mayo 2026 & Consolidación del \textit{firmware} de estación con lectura de sensores, modo seguro, \textit{cooldown}, decisión local y riego temporizado. \\
\hline
H11 & Junio 2026 & Montaje y prueba de la estación B, versión final del \textit{dashboard} y ajustes eléctricos del bloque de actuación. \\
\hline
H12 & 22 de junio de 2026 & Prueba funcional en Agrolab Móstoles con la estación B en entorno real de cultivo, grabación de vídeo, fotografías y recepción cualitativa. \\
\hline
H13 & Junio--julio 2026 & Cierre de memoria, anexos, resultados, discusión, conclusiones, entrega y preparación de defensa. \\
\hline
\end{longtable}

\normalsize

\section{Iteraciones, reajustes y problemas corregidos}

El desarrollo incluyó varios reajustes relevantes. Estos cambios no se consideran desviaciones negativas, sino decisiones técnicas necesarias para alcanzar una solución más coherente con el objetivo final del TFG.

\small
\renewcommand{\arraystretch}{1.12}
\rowcolors{2}{TableRowSoft}{white}

\begin{longtable}{|p{0.25\textwidth}|p{0.31\textwidth}|p{0.32\textwidth}|}
\caption{Iteraciones, reajustes y problemas corregidos durante el desarrollo del proyecto.}
\label{cuadro:iteraciones_cronograma}\\
\hline
\rowcolor{URJCRed}
\TVHead{Iteración o problema} & \TVHead{Motivo} & \TVHead{Resultado o corrección aplicada} \\
\hline
\endfirsthead

\hline
\rowcolor{URJCRed}
\TVHead{Iteración o problema} & \TVHead{Motivo} & \TVHead{Resultado o corrección aplicada} \\
\hline
\endhead

Reajuste desde arquitectura local hacia LTE/ThingSpeak & El enfoque inicial con Raspberry Pi, puerto serie y MQTT limitaba la validación de estaciones remotas en campo. & Se adoptó una arquitectura con estaciones autónomas, LTE, ThingSpeak y supervisión local posterior. \\
\hline
Decisión local de riego en el \textit{firmware} & Evitar depender de la nube, la Raspberry Pi o el \textit{dashboard} para accionar una bomba de agua. & La estación decide localmente y la supervisión queda separada del control crítico. \\
\hline
Problema de detección del Arduino compatible & La placa no aparecía inicialmente como puerto disponible por falta del controlador CH340G. & Se instaló el controlador correspondiente y se validó la carga de programas en Arduino IDE. \\
\hline
Estabilización de comunicación con A7670E-LASA & El módem requería pruebas de baudios, comandos AT, APN, registro LTE y publicación HTTP. & Se estabilizó la comunicación serie a \texttt{9600} baudios y se validó la publicación en ThingSpeak. \\
\hline
Separación entre publicación y control & ThingSpeak podía fallar temporalmente o no responder de forma inmediata. & La pérdida de publicación no bloquea la lógica local de riego ni genera una actuación insegura. \\
\hline
Evolución del \textit{dashboard} & La primera versión era suficiente para una demo, pero no para la supervisión final del sistema. & Se desarrolló una versión más completa con visualización, histórico, selección de estación y recursos gráficos. \\
\hline
Incorporación de estación B & Se quería demostrar que la arquitectura no dependía de una única estación física. & Se montó y probó una segunda estación, con configuración diferenciada y validación en Agrolab. \\
\hline
Ajustes en bloque de bomba y alimentación & La bomba introduce perturbaciones y exige una alimentación más robusta que la sensórica. & Se reforzó la documentación eléctrica y se incorporó un condensador de \(2200~\mu\mathrm{F}\) en la rama de la bomba. \\
\hline
Validación en Agrolab Móstoles & Era necesario probar el prototipo fuera del laboratorio y obtener evidencias reales de funcionamiento. & Se realizó una visita previa en abril y una prueba funcional el 22 de junio, con vídeo, fotografías y retroalimentación de usuarios. \\
\hline
\end{longtable}

\normalsize

\section{Dependencias entre tareas}

El desarrollo de TierraViva presentó dependencias claras entre bloques. Estas dependencias explican por qué algunas tareas tuvieron que resolverse antes de avanzar a fases posteriores.

\begin{itemize}
    \item La definición de la arquitectura final dependía de analizar previamente alternativas de comunicación, supervisión y actuación.
    \item La validación del \textit{firmware} dependía de estabilizar la lectura de sensores, el pinado, la alimentación y el comportamiento del relé.
    \item La publicación en ThingSpeak dependía de validar el A7670E-LASA, la tarjeta SIM, la cobertura, el APN, los comandos AT y la respuesta HTTP.
    \item El diseño de la base de datos y la API dependía de fijar previamente la estructura de campos ThingSpeak, nombres internos y códigos de estado.
    \item El \textit{dashboard} dependía de disponer de una API estable y de una estructura común para consultar estaciones diferenciadas.
    \item Las pruebas de bomba dependían de haber validado antes alimentación, masa común, relé, temporización y lógica de seguridad.
    \item La prueba en Agrolab Móstoles dependía de disponer de una estación funcional, comunicación validada, publicación de telemetría y evidencias suficientes para comprobar el flujo completo.
    \item La redacción final de resultados y conclusiones dependía de cerrar antes la validación funcional, los anexos técnicos y la trazabilidad entre requisitos y pruebas.
\end{itemize}

Estas dependencias muestran que el proyecto no se desarrolló como una suma de tareas independientes, sino como una integración progresiva. Cada bloque condicionó al siguiente y obligó a revisar decisiones a medida que el prototipo pasaba de una idea inicial a un sistema físico funcional.

\section{Redacción de la memoria}

La memoria se desarrolló en paralelo al avance técnico del proyecto, aunque la redacción intensiva comenzó en mayo de 2026. Primero se trabajaron los capítulos de introducción, estado del arte, requisitos y arquitectura; después los capítulos de implementación \textit{hardware}, \textit{firmware} y software; posteriormente se elaboraron los anexos técnicos, manuales, plan de pruebas, gestión de riesgos y trazabilidad; y finalmente se cerraron los capítulos de validación, resultados, discusión y conclusiones.

Esta forma de trabajo permitió que la documentación no se limitara a una descripción final del sistema, sino que recogiera también la evolución real del proyecto: decisiones descartadas, problemas encontrados, reajustes de arquitectura, validación por bloques, evidencias de funcionamiento y limitaciones detectadas.

\section{Conclusión del cronograma}

El cronograma refleja un desarrollo progresivo e iterativo desde finales de agosto de 2025 hasta julio de 2026. La primera parte del proyecto se centró en definir el problema, revisar tecnologías y plantear la arquitectura. La parte intermedia concentró el mayor trabajo técnico: sensores, Arduino, LTE, ThingSpeak, Raspberry Pi, SQLite, API, \textit{dashboard} y decisión local de riego. La fase final integró las estaciones, validó el sistema, generó evidencias, realizó la prueba en Agrolab Móstoles y cerró la memoria.

La planificación seguida muestra que TierraViva no fue un montaje directo ni una implementación lineal. El proyecto evolucionó mediante pruebas, correcciones y decisiones técnicas justificadas, hasta alcanzar un prototipo funcional de extremo a extremo.